\documentclass[border=0pt]{standalone}
\usepackage[RPvoltages]{circuitikz}
\usetikzlibrary{arrows,decorations.markings,arrows.meta}
\newcommand{\bushere}[3]{% length, text above, text below
    coordinate(tmp)
    % go up and do an edge down
    ++(0,#1) node[anchor=base]{#2} edge[line width=2.5pt] ++(0,{-1.88*#1})
    % edges do not move the current point, go down to position the node
    ++(0,{-1.88*#1}) node[below]{#3}
    % go back to where we started
    (tmp)
}

\newcommand{\redbushere}[3]{% length, text above, text below
%    coordinate(tmp)
    % go up and do an edge down
    ++(0.04,#1) node[anchor=base]{#2} edge[red,dashed, line width=1.5pt]  ++(0,{-1.88*#1})
    % edges do not move the current point, go down to position the node
    ++(0,{-1.88*#1}) node[below]{#3}
    % go back to where we started
%    (tmp)
}
\begin{document}\sf
%\ctikzset{sources/fill=gray!20, resistors/fill=gray!20}
\ctikzset{sources/fill=white, resistors/fill=white}
%\tikzset{myarr/.style={line width=0.74pt,-{Triangle[width=6pt, length=9pt]}}}
\tikzset{myarr/.style={line width=0.74pt,-{Triangle[width=6pt, length=9pt]}}}

\begin{circuitikz}[>=latex]% default line width
\ctikzset{bipoles/cuteswitch/thickness=0.4}
\draw [myarr](0,0) to [american voltage source,rotate=90,line width=0.4pt] 
 node[yshift=3ex]{G}(0,0)to (1.5,0);
\coordinate (B) at (1.5,0);
\coordinate (C) at (6.5,0);

\path (B)   \bushere{1}{}{} ;
%\path (B)   \redbushere{1}{}{} ;
 
%\draw [myarr,red](B)	++(0,0.5) 
%     to node[xshift=-0.5ex,yshift=1.3ex]{\color{black}$p_{lg}$} ++(1.5,0);
\draw [myarr](B)node[xshift=-4.2ex,yshift=1.3ex]{$p$}++(0,-0.6) 
     to [cosw,l_=CB] node[xshift=-18ex,yshift=1.3ex]{$p_1$} ++(5,0);
\path (C)  \bushere{1}{}{} ;
\path (C)  \redbushere{1}{}{} ;
\draw [myarr,red](C)	++(0,0.5) 
     to node[xshift=-0.5ex,yshift=1.3ex]{\color{black}$p_{l}$} ++(1.2,0);
\draw [myarr](C)++(0,-0.6) 
     to node[xshift=-0.5ex,yshift=1.3ex]{$p_2$} ++(2,0);

\end{circuitikz}
\end{document}